\begin{center}
  {\Large\ReportTitle\par}
  \vspace{0.15cm}
  {\large\StudentName\par}
  \vspace{0.15cm}
  {\large Degree \DegreeName\par}
  \vspace{0.15cm}
  {\large\InstitutionName\par}
  \vspace{0.15cm}
  {\large\CopyrightYear\par}
\end{center}

\vspace{0.8cm}
\phantomsection
\addcontentsline{toc}{chapter}{Abstract}
\noindent{\huge\bfseries Abstract\par}
\vspace{0.8cm}

Enterprise question-answering systems must retrieve organization-specific
knowledge, follow business workflow dependencies, and provide evidence that users can
trace back. A static Retrieval-Augmented Generation (RAG) pipeline applies the
same retrieval strategy to every request, even though business questions range
from direct definitions to multi-document and tool-assisted investigations.
This project designs and implements \SystemName, an Adaptive RAG system for
business workflow question answering. The system predicts query complexity
and supports four execution levels: direct generation, single-step retrieval,
decomposed multi-query retrieval, and a bounded agentic route.

The implementation combines a React web-based RAG studio, a FastAPI service,
a PostgreSQL relational database, pgVector retrieval, and an LLM model farm.
Its knowledge pipeline supports
document ingestion, metadata extraction, deduplication, versioned indexes,
token-aware structure-preserving chunking, embedding, and hybrid BM25/dense
retrieval. Query-time processing adds conversation-aware reformulation,
optional reranking, context budgeting, grounded prompts, citation validation,
streaming generation, and full-fidelity execution traces. 
Evaluation integrates the enterprise-oriented WixQA benchmark with lexical, 
retrieval, judge-based, latency, token, and cost measurements. 
RAGXplain converts metric outputs into diagnostic findings and
configuration-level recommendations.

Both DistilBERT and T5-small complexity classifiers achieved 98.75\% accuracy and 98.74\%
macro F1 on a 481 validation records. In the primary 20-case, 
100-document WixQA configuration experiment,
Adaptive routing with DistilBERT ranked first with a quality score of 0.5449.
RAGXplain consistently identified retrieval coverage and grounding as the
principal improvement targets. However, it remains limited by the compatible sample size 
and configuration differences among the compared systems.

\noindent\textbf{Keywords:} Adaptive RAG, business workflow question
answering, query complexity, WixQA, RAGXplain, explainable
evaluation, large language models.
